\documentclass[11pt]{article}

\usepackage[T1]{fontenc}
\usepackage[margin=1in]{geometry}
\usepackage{hyperref}
\usepackage{xcolor}
\usepackage{listings}
\usepackage{tabularx}

\definecolor{codegray}{RGB}{248,248,248}
\lstset{
  basicstyle=\small\ttfamily,
  backgroundcolor=\color{codegray},
  frame=single,
  breaklines=true,
  columns=fullflexible,
  keepspaces=true,
  showstringspaces=false
}

\title{Displaying Uninitialized Value Flows in ESBMC Counterexamples}
\author{ESBMC}
\date{}

\begin{document}

\maketitle

\section{Overview}

ESBMC normally slices assignments that do not affect the violated property.
This keeps the verification formula and counterexample small, but it can hide
program state that is useful during counterexample replay.  For example, the
following assignments do not affect the failing assertion:

\begin{lstlisting}[language=C]
unsigned initialized;
unsigned counter = initialized;
counter++;
\end{lstlisting}

The \texttt{--show-uninitialized} option retains the complete value flow that
starts from an uninitialized local variable.  The counterexample shows the
concrete value selected for the uninitialized variable, the copy into
\texttt{counter}, and the subsequent increment.  The option is disabled by
default, so existing counterexamples and slicing behaviour remain unchanged.

\section{Example}

Save the following program as \texttt{example.c}:

\begin{lstlisting}[language=C]
int main(void)
{
  unsigned initialized;
  unsigned counter = initialized;
  counter++;

  unsigned reset = 0;
  __ESBMC_init_object(&reset);
  unsigned reset_counter = reset;
  reset_counter++;

  __ESBMC_assert(0, "expected failure");
  return 0;
}
\end{lstlisting}

This example contains two uninitialized-value sources:

\begin{itemize}
  \item \texttt{initialized} is declared without an initializer.
  \item \texttt{\_\_ESBMC\_init\_object(\&reset)} replaces the value of
        \texttt{reset} with an unconstrained value.
\end{itemize}

Assignments derived from either source are retained even though they do not
affect the final assertion.

\section{Command}

Run ESBMC with the new option:

\begin{lstlisting}[language=bash]
esbmc example.c --incremental-bmc --show-uninitialized
\end{lstlisting}

When running ESBMC directly from a build tree, use:

\begin{lstlisting}[language=bash]
./build/src/esbmc/esbmc example.c \
  --incremental-bmc --show-uninitialized
\end{lstlisting}

No additional slicing option, including \texttt{--slice-assumes}, is needed.

\section{Counterexample Output}

A representative counterexample is shown below.  The exact concrete value is
solver-dependent; in this run, the solver selected zero for both unconstrained
values.

\begin{lstlisting}
[Counterexample]

State 1 file example.c line 3 function main thread 0
----------------------------------------------------
  initialized = 0 (00000000 00000000 00000000 00000000) (uninitialized)

State 2 file example.c line 4 function main thread 0
----------------------------------------------------
  counter = 0 (00000000 00000000 00000000 00000000)

State 3 file example.c line 5 function main thread 0
----------------------------------------------------
  counter = 1 (00000000 00000000 00000000 00000001)

State 4 file example.c line 8 function main thread 0
----------------------------------------------------
  reset = 0 (00000000 00000000 00000000 00000000) (uninitialized)

State 5 file example.c line 9 function main thread 0
----------------------------------------------------
  reset_counter = 0 (00000000 00000000 00000000 00000000)

State 6 file example.c line 10 function main thread 0
----------------------------------------------------
  reset_counter = 1 (00000000 00000000 00000000 00000001)

State 7 file example.c line 12 function main thread 0
----------------------------------------------------
Violated property:
  expected failure
  0

VERIFICATION FAILED
\end{lstlisting}

Only the assignments that introduce unconstrained values carry the
\texttt{(uninitialized)} marker.  Derived variables such as \texttt{counter}
and \texttt{reset\_counter} are initialized variables, so they are displayed
without the marker.  Their assignments remain visible because their values
depend on an uninitialized source.

\section{Validated Object Types}

The \texttt{\_\_ESBMC\_init\_object} counterexample support has been tested
with scalar and aggregate C objects.  Table~\ref{tab:init-object-support}
summarizes the observed results.

\begin{table}[ht]
\centering
\begin{tabularx}{\textwidth}{|l|X|l|}
\hline
Object type & Example & Status \\
\hline
Integer variable & \texttt{int value} & Supported \\
\hline
Array & \texttt{int values[3]} & Supported \\
\hline
Structure & \texttt{struct pair} & Supported \\
\hline
Array of structures & \texttt{struct item values[2]} & Supported \\
\hline
Nested object & Structure containing another structure or an array &
Supported with Boolector; partially concretized with Z3 \\
\hline
Mixed field types & Structure containing \texttt{int}, \texttt{double}, and
\texttt{char} & Supported \\
\hline
\end{tabularx}
\caption{Validated \texttt{\_\_ESBMC\_init\_object} object types.}
\label{tab:init-object-support}
\end{table}

Both array call forms were tested:

\begin{lstlisting}[language=C]
__ESBMC_init_object(values);
__ESBMC_init_object(&values);
\end{lstlisting}

In both cases, the array is displayed as an uninitialized-value source and
reads derived from its elements remain in the counterexample.

\subsection{Validation Example}

The following small structure demonstrates the feature:

\begin{lstlisting}[language=C]
struct data
{
  char tag;
  int count;
};

int main(void)
{
  struct data value = {0, 0};
  __ESBMC_init_object(&value);
  __ESBMC_assume(value.tag == 'A');
  __ESBMC_assume(value.count == 3);
  __ESBMC_assert(0, "expected failure");
}
\end{lstlisting}

Run it with:

\begin{lstlisting}[language=bash]
esbmc example.c --incremental-bmc --show-uninitialized
\end{lstlisting}

The counterexample contains the complete object value:

\begin{lstlisting}
[Counterexample]

State 1 file example.c function main thread 0
----------------------------------------------------
  value = { .tag=65, .count=3 } (uninitialized)

VERIFICATION FAILED
\end{lstlisting}

The decimal value 65 is the character code for \texttt{'A'}.

\section{What Is Considered Uninitialized}

A local declaration without an initializer is an uninitialized source:

\begin{lstlisting}[language=C]
int value;
\end{lstlisting}

An explicit call to \texttt{\_\_ESBMC\_init\_object} is also an uninitialized
source:

\begin{lstlisting}[language=C]
int value = 0;
__ESBMC_init_object(&value);
\end{lstlisting}

An explicitly initialized nondeterministic variable is \emph{not} classified
as uninitialized:

\begin{lstlisting}[language=C]
int value = nondet_int();
\end{lstlisting}

The last assignment may still appear when it is relevant to a violated
property, but it will not carry the \texttt{(uninitialized)} marker and does
not become a root for this value-flow display feature.

\section{Behaviour Without the Option}

Running the same example without \texttt{--show-uninitialized}:

\begin{lstlisting}[language=bash]
esbmc example.c --incremental-bmc
\end{lstlisting}

uses the normal slicing rules.  Since the assignments to \texttt{initialized},
\texttt{counter}, \texttt{reset}, and \texttt{reset\_counter} do not affect the
failed assertion, they are omitted from the counterexample.  The output still
ends with:

\begin{lstlisting}
Violated property:
  expected failure
  0

VERIFICATION FAILED
\end{lstlisting}

\end{document}
